Я изложу основные идеи метода Flow Matching и приведу доказательства некоторых утверждений
Пусть:
\begin{itemize}
    \item $p_0(x)$ — простое распределение (например, $\mathcal N(0,I)$),
    \item $p_1(x)$ — целевое распределение данных.
\end{itemize}
Задача: построить диффеоморфизм
\[
\psi_1 : \mathbb R^d \to \mathbb R^d,
\]
такой что
\[
p_1(x) = [\psi_1]_\# p_0 (x) = p_0(\psi_1^{-1}(x)) |\partial_x \psi_1^{-1}(x)|
\]
Т.е. $X_1 = \psi_1(X_0) \sim p_1 \iff X_0 \sim p_0$. Вместо того, чтобы строить $\psi_1$ напрямую, рассмотрим векторное поле
\[
\frac{d}{dt}\psi_t(x) = v_t(\psi_t(x)), \quad t \in [0,1].
\]
Из которого отображение $\psi_1$ восстанавливается интегрированием
\[
\psi_1(x) = p_0(x) + \int_0^{1} v_t(\psi_t(x)) dt
\]

TODO: добавить теорему о существовании и единственности такого векторного поля из 2412.06264v1
\begin{theorem}

\end{theorem}

\begin{proposition}
    Если $\psi_t$ удовлетворяет
    \[
    \frac{d}{dt}\psi_t(x) = v_t(\psi_t(x)),
    \]
    то плотность $p_t = [\psi_t]_\# p_0 (x)$ удовлетворяет уравнению непрерывности:
    \[
    \partial_t p_t(x) + \nabla \cdot \big(p_t(x) v_t(x)\big) = 0.
    \]
\end{proposition}
Тогда, если мы найдём правильное $v_t$, то $\psi_t$ \textit{хорошим образом} перенесёт $p_0$ в $p_1$
\begin{proof}
    Поскольку $\rho_t (y)$ определено как 
    \[
    \rho_t (y) = (\psi_t)_\# \rho_0 (y) = \rho_0(\psi_t^{-1}(y)) | \Det \partial_y \psi_t^{-1}(y)|
    \]
    то 
    \[
    \int \varphi(y)\,\rho_t(y)\,dy = \int \varphi(\psi_t(x))\,\rho_0(x)\,dx \quad \forall \varphi \in C_c^\infty.
    \]
    Продифференцируем обе части по $t$:
    \[
    \frac{d}{dt} \int \varphi(\psi_t(x))\,\rho_0(x)\,dx = \int \frac{d}{dt}\varphi(\psi_t(x))\,\rho_0(x)\,dx
    \]
    \[
    \frac{d}{dt}\varphi(\psi_t(x)) = \nabla \varphi(\psi_t(x)) \cdot \dot{\psi}_t(x).
    \]
    И т.к.
    \[
    \dot{\psi}_t(x) = u_t(\psi_t(x)),
    \]
    получим
    \[
    \frac{d}{dt} \int \varphi(\psi_t(x))\,\rho_0(x)\,dx = \int \nabla \varphi(\psi_t(x)) \cdot u_t(\psi_t(x))\,\rho_0(x)\,dx
    \]
    Снова, по определению $\rho_t (y)$
    \[
    \int \nabla \varphi(\psi_t(x)) \cdot u_t(\psi_t(x))\,\rho_0(x)\,dx = \int \nabla \varphi(y) \cdot u_t(y)\,\rho_t(y)\,dy
    \]
    Интегрированием по частям получим
    \[
    \int \nabla \varphi \cdot u_t \rho_t = - \int \varphi \, \nabla \cdot (\rho_t u_t)
    \]
    Тогда
    \[
    \frac{d}{dt} \int \varphi \rho_t + \int \varphi \, \nabla \cdot (\rho_t u_t) = 0
    \]
    И посольку это выполнено для всех $\varphi$, то
    \[
    \partial_t \rho_t + \nabla \cdot (\rho_t u_t) = 0
    \]
    в слабом смысле
\end{proof}

\subsection{Ключевая идея Flow Matching}
\hfill \break
\hfill \break
Вместо обучения напрямую распределения $q_\theta(x) \approx p_1(x)$ через KL-дивергенцию,
Flow Matching предлагает:
\begin{enumerate}
    \item Задать интерполирующее распределение $p_t$ между $p_0$ и $p_1$
    \item Вычислить соответствующее \textit{истинное} поле скоростей $u_t(x)$.
    \item Обучать модель $v_\theta$ по регрессии:
\end{enumerate}
\[
\mathcal L_{\text{FM}} =
\mathbb E_{t \sim U[0,1],\, x \sim p_t}
\| v_t^\theta(x) - u_t(x) \|^2.
\]
Проблема лишь в том, что \textit{истинное} поле $u_t(x)$ нам обычно неизвестно.

\subsection{Условная формулировка}
\hfill \break
\hfill \break
Тогда, вместо маргинального полея $u_t(x)$ введем условное поле
\[
u_t(x \mid z),
\]
где $z$ — латентная случайная величина (например, пара $z = (x_0, x_1) \sim \pi_{0, 1} $). Тогда условный лосс:
\[
\mathcal L_{\text{cond}} =
\mathbb E_{t,z,x_t \sim p_{1|z}(\cdot \mid z)}
\| v_t^\theta(x_t) - u_t(x_t \mid z) \|^2.
\]
И это легко сэмплировать:
\begin{enumerate}
    \item берём $z$,
    \item генерируем $x_t \sim p_{1 | z}$,
    \item знаем $u_t(x_t \mid z)$ аналитически
\end{enumerate}
В таком случае истинное поле выражается как
    \[
    u_t(x)
    =
    \mathbb{E}_{z \sim \pi} \big[
    u_t(X_t \mid z) \mid X_t = x
    \big].
    \]

\subsection{Пример}
\hfill \break
\hfill \break
Простейший случай с двумя гауссианами

\begin{align*}
    p_0(x) &= \mathcal N(\mu_0, \sigma_0) \\
    p_1(x) &= \mathcal N(\mu_1, \sigma_1) \\
    \psi_{t \mid 1}(x \mid x_1) &= t x_1 + (1 - t) x
\end{align*}

Тогда условное поле скоростей:
    \[
    u_{t \mid 1}(x \mid x_1) = \dot{\psi}_{t \mid 1} \circ \psi_{t \mid 1}^{-1}(x \mid x_1) = \frac{x_1 - x}{1 - t}
    \]
И условный probability path:
    \[
    p_{t \mid 1} (x \mid x_1) = p_{0 \mid 1} (\psi_{t \mid 1}^{-1}(x \mid x_1)) | \partial_x\psi_{t \mid 1}^{-1}(x \mid x_1) | = 
    \]
    \[
    = \frac{1}{\sqrt{2 \pi} \sigma} e^{-\frac{1}{2} \Big( \frac{\frac{x - t x_1}{1 - t} - \mu_0}{\sigma_0} \Big)^2} \cdot \frac{1}{1 - t} = 
    \mathcal N(t x_1 + (1 - t)\mu_0, (1 - t)\sigma_0)
    \]
Тогда marginal probability path:
    \[
    p_t(x) = \int p_{t \mid 1}(x | x_1) p_1(x_1) dx_1 = 
    \]
    \[
    = \int \mathcal N(t x_1 + (1 - t)\mu_0, (1 - t)\sigma_0) \> \mathcal N(\mu_1, \sigma_1) dx_1 =
    \]
    \[
    = \mathcal{N}(t \mu_1 + (1 - t) \mu_0, t^2 \sigma_1 + (1 - t) \sigma_0)
    \]
Отсюда по формуле Байеса 
    \[
    p_{1 \mid t}(x_1 \mid x) = \frac{p_{t \mid 1}(x \mid x_1) p_1(x_1)}{p_t(x)} = \mathcal N(\mu_{1 | t}(x), \sigma_{1|t}(x))
    \]
где
    \[
    \mu_{1 \mid t}(x) = \mu_1 + \frac{t \sigma_1 (x - t\mu_1 - (1 - t)\mu_0)}{t^2 \sigma_1 + (1 - t)\sigma_0}
    \]
    \[
    \sigma_{1 \mid t}(x) = \frac{\sigma_1 (1 - t) \sigma_0}{t^2 \sigma_1 + (1 - t)\sigma_0}
    \]
Итого, marginal velocity field:
    \[
    u_t(x) = \int u_{t \mid 1}(x \mid x_1) p_{1 \mid t}(x_1 \mid x) dx_1 = 
    \]
    \[
    = \mathbb{E}\Big[\frac{X_1 - X_t}{1 - t} \mid X_t = x\Big] = \frac{1}{1 - t}\mathbb{E}[X_1 \mid X_t = x] - \frac{x}{1 - t} =
    \]
    \[
    = \frac{\mu_{1 | t}(x) - x}{1 - t}
    \]

\subsection{Теорема о совпадении лоссов}
% \hfill \break
\begin{theorem}{\cite{DINGO}}

Градиент условного лосса совпадает с градиентом маргинального лосса:
\[
\nabla_\theta \mathcal L_{\text{cond}}
=
\nabla_\theta \mathcal L_{\text{marg}},
\]
где
\[
\mathcal L_{\text{cond}} =
\mathbb E_{t,Z \sim \pi, X_t \sim p_{1 | Z}}
\| v_t^\theta(X_t) - u_t(X_t|Z) \|^2.
\]

\[
\mathcal L_{\text{marg}}
=
\mathbb E_{t, X_t \sim p_t}
\| v_t^\theta(X_t) - u_t(X_t) \|^2
\]
\end{theorem}

\begin{proof}{\cite{DINGO}}
\begin{align*}
        \nabla_\theta \mathcal L_{\text{marg}} 
        & = \nabla_\theta \mathbb{E}_{t,X_t} \| v_t^\theta(X_t) - u_t(X_t) \|^2 \\
        & = \mathbb{E}_{t,X_t} \nabla_\theta \| v_t^\theta(X_t) - u_t(X_t) \|^2 \\
        & = \mathbb{E}_{t,X_t} \textcolor{blue}{\nabla_\theta v_t^\theta(X_t) \nabla_{v_t^\theta}} \| v_t^\theta(X_t) - u_t(X_t) \|^2 \\
        & = \mathbb{E}_{t,X_t} \nabla_\theta v_t^\theta(X_t) \nabla_{v_t^\theta} \| v_t^\theta(X_t) - \textcolor{blue}{\mathbb{E}_{Z \sim p_{Z | t}} [ u_t(X_t \mid Z)]} \|^2 \\
        & = \mathbb{E}_{t,X_t} \textcolor{blue}{\mathbb{E}_{Z \sim p_{Z | t}}} \nabla_\theta v_t^\theta(X_t) \nabla_{v_t^\theta} \| v_t^\theta(X_t) - u_t(X_t \mid Z) \|^2 \\
        & = \mathbb{E}_{t,X_t} \textcolor{blue}{\mathbb{E}_{Z \sim p_{Z | t}}} \nabla_\theta \| v_t^\theta(X_t) - u_t(X_t \mid Z) \|^2 \\
        & = \mathbb{E}_{t, \textcolor{blue}{X_t \sim p_{t|Z}, Z \sim \pi}} \nabla_\theta \| v_t^\theta(X_t) - u_t(X_t \mid Z) \|^2 \\
        & = \nabla_\theta \mathbb{E}_{t, X_t \sim p_{t|Z}, Z \sim \pi} \| v_t^\theta(X_t) - u_t(X_t \mid Z) \|^2 \\
        & = \nabla_\theta \mathcal{L}_{\text{cond}} 
\end{align*}
\end{proof}
